\begin{table}[htbp]
\centering
\caption{RQ1a protocol sensitivity. Identical encoder, budget and seed; only the split changes. Random is the common default in circuit representation-learning evaluations, recipe-disjoint is OpenABC-D Variant~1, design-disjoint is Variant~2 and is what this thesis reports.}
\label{tab:rq1a_protocol}
\small
\begin{tabular}{lrrrr}
\toprule
\textbf{Protocol} & \textbf{RMSE} & \textbf{$R^2$} & \textbf{Spearman} & \textbf{Inflation vs. design-disjoint} \\
\midrule
\textcolor{red}{{[TODO/FAKE] Random (per-row)}} & \textcolor{red}{999999} & \textcolor{red}{-999999} & \textcolor{red}{-999999} & \textcolor{red}{999999} \\
\textcolor{red}{{[TODO/FAKE] Recipe-disjoint}} & \textcolor{red}{999999} & \textcolor{red}{-999999} & \textcolor{red}{-999999} & \textcolor{red}{999999} \\
Design-disjoint (ours) & 0.043 & 0.343 & 0.185 & 0.000 \\
\bottomrule
\end{tabular}
\\[4pt]\begin{minipage}{\linewidth}\footnotesize\emph{\textcolor{red}{Rows in red are FABRICATED PLACEHOLDERS, not measurements: every number in them has been replaced by the sentinel $\pm$999999.} Rows marked [TODO/FAKE] are invented. IN FLIGHT: the --split\_by random and --split\_by recipe training runs are on the cluster now (the design-disjoint run already exists). Identical encoder, budget and seed; only the split changes. When they land, run src/test.py for each, delete this block and rerun analysis.make\_all.}\end{minipage}
\end{table}
